\documentclass[a4paper,oneside]{report}

\renewcommand{\chaptername}{Week}

\usepackage[margin=2.5cm]{geometry}
\usepackage{parskip}

\usepackage{amsmath}
\usepackage{amssymb}

\usepackage{pgfplots}

\usepackage{listings}
\lstset{
	basicstyle=\ttfamily,
	numbers=left,
	escapechar=\%,
}

\usepackage{tikz}
\usetikzlibrary{arrows, positioning, shapes}
\tikzstyle{goal} = [draw, ellipse, minimum height=1cm, minimum width=2cm]
\tikzstyle{process} = [draw, rectangle, rounded corners, minimum height=1cm,
	minimum width=1cm]
\tikzstyle{decision} = [draw, diamond, minimum height=1cm, minimum width=2cm]
\tikzstyle{line} = [draw, -latex']

\begin{document}

\begin{titlepage}
	\raggedleft
	\vspace*{1cm}
	\Huge{Computational Mathematics}\\
	\vspace{0.5cm}
	\Large{ECM1416}
\end{titlepage}

\null
\thispagestyle{empty}
\addtocounter{page}{-2}
\newpage

\tableofcontents

\chapter{Vectors and Introduction to Matrices}

\section{Notes}

\subsection{Asynchronous Videos}

We have previously seen that graphs can be formed using a set of axis, for
example the 2 dimensional coordinate system requires a pair of perpendicular
axis which share a single unit of length across them. Typically the horizontal
axis is labelled \( x \) and the vertical axis \( y \). The point where these
axis meet is called the origin and is labelled \( O \). All other coordinates
are defined by an ordered pair of numbers, a tuple whose first number defines
how far to move along the \( x \)-axis and the second how far to move along
the \( y \)-axis. You can think of the axis as being a pair of number lines,
with the origin having coordinates (0, 0) and all other coordinates following
from that. The following graph displays some coordinates:

\begin{figure}[h]
	\centering
	\begin{tikzpicture}
		\begin{axis}[axis y line=center, axis x line=center,
			xtick={-5,-4,...,5}, ytick={-5,-4,...,5},
			xmin=-5.5, xmax=5.5, ymin=-5.5, ymax=5.5,
			xlabel=\( x \), ylabel=\( y \)]
			\addplot[mark=*] coordinates {(0, 0)};
			\addplot[mark=*] coordinates {(3, 3)};
			\addplot[mark=*] coordinates {(-3, 4)};
			\addplot[mark=*] coordinates {(-3, -2)};
			\addplot[mark=*] coordinates {(4, -2)};
		\end{axis}
	\end{tikzpicture}

	The points \((0, 0), (3, 3), (-3, 4), (-3, -2), \text{ and }
	(4, -2)\) are shown here.
\end{figure}

This coordinate system can be defined more thoroughly to display coordinates
whose components are real numbers, creating a real coordinate space. The above
is an example of the real coordinate space of dimension 2, also known as
\( \mathbb{R}^2 \). A real coordinate space of dimension \( n \) is written as
\( \mathbb{R}^n \), and the coordinates on this space are an ordered tuple
are given by \( (x_1, x_2, x_3, \dots, x_n) \).

Let the coordinate \( P = (7, -4) \) and the coordinate \( Q = (-3, 6) \). We
can find a vector \( \overrightarrow{PQ} \) by doing \( P - Q \) element-wise,
that is \( 7 - -3 = 10 \) and \( -4 - 6 = -10 \). The vector
\( \overrightarrow{PQ} = (10, -10) \) where the first number defines how far
to move along the \( x \)-axis and the second how far to move along the
\( y \)-axis. This is similar to the syntax used for coordinates, however
unlike coordinates this is not limited to the axis. You can think of a vector
as being instructions of how to move between a particular set of points, or
alternatively how to move in a well defined way. Some important vectors are:

\begin{itemize}
	\item{Unit vectors: these are vectors whose norm is equivalent to
		1 of the unit that is used by the axis on which they lie. For
		a given vector \( \overrightarrow{u} \) its unit vector is
		written as \( \hat{u} \). Provided that the original vector is
		not the zero vector, this unit vector will point in the same
		direction as the original vector.}
	\item{Basis vectors: in \( \mathbb{R}^2 \), these are typically (1, 0)
		and (0, 1). You may think of these as movements along the
		\( x \)- and \( y \)-axis respectively, and they form the
		basis for all the other vectors in \( \mathbb{R}^2 \). They
		don't have to be these vectors however; any 2 vectors \( a \)
		and \( b \) can form the basis vectors of \( \mathbb{R}^2 \)
		provided they are not paralell or zero vectors. All the other
		vectors in \( \mathbb{R}^2 \) can then be defined as multiples
		of these vectors and so every point in \( \mathbb{R}^2 \) can
		be reached.}
	\item{Zero vectors: these are vectors whose norm is 0, and so all its
		components are also 0.}
	\item{\( \overrightarrow{OA} \): recall that \( O \) is used to refer
		to the origin of a given coordinate space. Let \( A \) be a
		coordinate in the coordinate space. The vector
		\( \overrightarrow{OA} \) therefore refers to the vector that
		moves from the origin to the point \( A \).}
\end{itemize}

Generally speaking, vectors are often referred to with an arrow overhead,
for example \( \overrightarrow{a} = (1, 3, -4) \). The components/elements of
a vector are referred to using subscripts, for example \( a_2 = 3\).

Recall unit vectors, these are vectors whose norm is 1. The norm of a vector
refers to its length/magnitude. To find the norm of a vector
\( \overrightarrow{v} \), we do the following:

\[ |\overrightarrow{v}| = \sqrt{v_1^2 + v_2^2 + \dots + v_n^2} \]

For example, the norm of (3, 4) is \( \sqrt{3^2 + 4^2} = 5 \). We can use this
result to find the unit vector by doing the following:

\[ \hat{u} = \frac{\overrightarrow{u}}{|\overrightarrow{u}|} \]

Using the example from above, we get the unit vector \( (\frac{3}{5},
\frac{4}{5}) \). We can confirm this is in fact a unit vector by finding its
norm:

\[ \sqrt{ \frac{3}{5}^2 + \frac{4}{5}^2 } = 1 \]

Technically, there are many kinds of norm. This specific kind is called the
\textit{Euclidian norm}, and generally if a norm is mentioned without some
prefix this is the norm that is being referred to.

Some important properties of a norm:

\begin{itemize}
	\item{The norm of a vector \( \overrightarrow{v} \geq 0 \). This
		should be obvious from the fact that we square a set of values
		and then take the square root of these values.}
	\item{\( |k\overrightarrow{v}| = |k||\overrightarrow{v}| \) where
		\( k \) is some constant scalar quantity.}
	\item{\( |\overrightarrow{u} + \overrightarrow{v}| \leq
		|\overrightarrow{u}| + |\overrightarrow{v}| \): this is
		the triangle inequality. Essentially, one side of a triangle
		will always be shorter than or equal to the other 2 sides
		added together.}
	\item{\( |\overrightarrow{u} - \overrightarrow{v}| \geq
		|\overrightarrow{u}| - |\overrightarrow{v}| \) this is
		the reverse triangle inequality.}
\end{itemize}

Still need to add notes for other videos, this is just the first one.

\subsection{Synchronous Lectures}

\section{Workshops}

Workshop exercises for this week go here.

\section{Exercises}

Solution to exercise sheets set this week go here.

\end{document}
